\documentclass[12pt]{article}

\usepackage[utf8]{inputenc}
\def\magyarOptions{defaults=hu-min}
\usepackage[magyar]{babel}

\usepackage{ragged2e}
\usepackage{t1enc}
\usepackage[a4paper, margin=2.5cm]{geometry}

\title{Bus4U - UI}
\author{Portik Szabolcs}
\date{}

\begin{document}
  \maketitle

  Bus4U - User Interface

  \begin{abstract}
    \justifying
    A webalkalmazás a tömegközlekedés megkönnyítését szolgálja. Lehetőséget biztosít jegyvásárlásra egy megadott járatra, illetve megtekinthetőek a járatok időpontjai városokra és megállókra lebontva. Ezen kívül megtekinthetőek a buszmegállók lokációi városok, illetve útvonalak szerint szűrve. Egy live map funkcióban valós időben követhetőek az autóbuszok jelenlegi helyzete, hogy a felhasználó előre informálódni tudjon a menetrend változásáról. Utóbbihoz, illetve a jegykezeléshez szükség van egy alkalmazásra, ami a buszokon található POS terminálon fut. Ez biztosítja a jegyek érvényesítését, illetve az autóbuszok helyzetének megosztását.
    \newline
    Ezen felül a webalkalmazás biztosít egy adminisztrációs felületet a busztársaságok számára. Ebben lehetőség nyílik elsősorban a menetrendek kezelésére, megállók és útvonalak feltöltésére. Itt lehet kezelni az alkalmazottak információit, és az autóbuszok különböző adatait.
    \newline
    Ezen dolgozat a UI részével foglalkozik.
  \end{abstract}

\end{document}